\section{无标签中子电流档案的可追溯观测与证据约束解释}

本研究将两类性质不同的资源置于同一诊断辅助流程中：一类是核电厂运行环境中积累的中子电流历史记录，用于形成可复核的信号现象；另一类是核仪表与控制技术资料，用于提供可定位的文献证据。二者都不是故障标签库。前者不能仅凭波形确认设备故障，后者也不能脱离测量现象直接解释某段原始信号。本章据此明确研究对象、两侧任务及其受控连接方式。

\subsection{研究对象、数据结构与披露边界}

测量资源来自同一核电厂采集周期内的中子电流历史档案。档案包含 42 条多通道 CSV 记录，来自四套仪表组。每条记录由时间列和 7 个中子电流通道组成；典型记录约含 $6\times10^{5}$ 个样本，采样间隔约为 1 s，并覆盖多日窗口。组内通道具有相对布置关系，因此其顺序不是任意的特征排列。原始导出文件可能含有局部时间戳紊乱或非单调片段，时间顺序校正因而是后续分析的前提。

档案不随每段时序提供完整的运行工况、维护历史和故障确认记录。反应堆型号、厂站身份、具体仪表系统名称及其他可识别信息亦暂不公开；本章只保留理解数据组织、方法适用范围和结果解释所必需的信息。本文据此关注可被复核的时间变化、通道间关系和信号完整性线索，而不依据记录推断反应堆型号或直接重建反应堆物理状态。

\subsection{测量侧：从原始记录到 observation packet}

测量侧的产物不是故障分类结果，而是一份 \textbf{observation packet}：它说明对哪一条记录、以哪些确定性规则、观察到了哪些数据质量和信号现象。确定性脚本将结构化观测写入记事板；其中稳定字段直接形成供后续检索使用的 compact alert summary，测量侧 agent 则只读取这些已保存结果，组织面向工程师阅读的 condition report。二者均不构成对物理机理、设备状态或根因的确认。

\begin{figure*}[!t]
  \centering
  \begin{minipage}[t]{0.263\textwidth}
    \centering\textbf{(a)}\par\vspace{1mm}
    \includegraphics[width=\linewidth]{fig03a_spike_rule_example}
  \end{minipage}\hfill
  \begin{minipage}[t]{0.341\textwidth}
    \centering\textbf{(b)}\par\vspace{1mm}
    \includegraphics[width=\linewidth]{fig03b_pearson_example}
  \end{minipage}\hfill
  \begin{minipage}[t]{0.296\textwidth}
    \centering\textbf{(c)}\par\vspace{1mm}
    \includegraphics[width=\linewidth]{fig03c_lag_example}
  \end{minipage}
  \caption{确定性 probes 所描述的统计现象示例：(a) 样本 2600--2720 的局部稳健偏离规则，纵轴为相应通道记录的中子电流值；(b) 通道间 Pearson 共变；(c) 相关--滞后扫描。它们用于说明 observation packet 中可保存的描述性证据，不构成已确认故障或因果方向。}
  \label{fig:probe-examples}
\end{figure*}

\subsubsection{初始化、数据健康与分析单元}

每次分析以一条记录为单位。初始化脚本解析时间列、按时间排序，并生成唯一的 active record；同时记录样本数、名义采样间隔、缺失情况和不活跃通道等数据健康信息。后续分析不再接受自由指定的原始文件路径，而是从该分析单元读取 active record。这样可以使一组观测始终对应同一输入记录，并保留输入、参数、脚本版本和派生产物之间的关联。

这一安排并不要求忽略数据结构或通道相对位置；相反，通道标识和组内相对关系是组织观察所必需的先验。其限制在于，测量侧不把这些结构信息延伸为预设故障知识，也不把任一统计量直接翻译为设备故障类别。

\subsubsection{预先定义的确定性分析工具}

初始化完成后，脚本仅执行预先注册的 probes，覆盖数据健康、形态与共变、事件与区段、同步与滞后、极值定位和截面快照等观察。各工具将输入范围、参数与结构化输出写入 observation packet；图~\ref{fig:probe-examples} 展示其中三类统计现象。形态和共变扫描提供的是信号关系的初步描述。例如，较高的两两相关性可以提示共享变化、可能的冗余或需要结合通道位置进一步核查的情形，但它本身不证明仪表关系或物理机制。相同地，相关—滞后扫描只区分近同步与时间错位的统计模式，不赋予因果方向。

为使事件识别可重复，当前实现对每个通道采用长度为 31 个样本的中心滑动窗口；记录边界处使用可获得的窗口样本。令 $W_t$ 为时刻 $t$ 的窗口，$x_{t,c}$ 为通道 $c$ 的电流值，则局部基线、点态绝对残差及其稳健局部尺度定义为

\[
\begin{aligned}
m_{t,c} &= \operatorname{median}_{\tau\in W_t}x_{\tau,c},\\
r_{t,c} &= \left|x_{t,c}-m_{t,c}\right|,\\
d_{t,c} &= \operatorname{median}_{\tau\in W_t}r_{\tau,c},\\
R_c &= \max_t x_{t,c}-\min_t x_{t,c}.
\end{aligned}
\]

这类以中位数和绝对偏差抑制局部异常影响的稳健估计思想，可追溯至相关基础研究 \citep{hampel1974influence,rousseeuw1993alternatives}。对于非恒定通道，突刺候选定义为同时满足局部与全程幅度条件的点：

\[
I^{\mathrm{spike}}_{t,c}=
\mathbb{I}\!\left(r_{t,c}>5d_{t,c}\right)
\mathbb{I}\!\left(r_{t,c}>0.05R_c\right).
\]

连续的阳性点合并为一个短时偏离事件；长度不超过两个采样点的连续零值段另行记录。该规则同时响应向上和向下的短时偏离，因而此处的“突刺”仅指检测事件，不预设其物理方向或故障含义。两通道之间以 Pearson 相关系数描述共变，当前实现将 $r\geq0.99$ 标记为近冗余，将 $0.85\leq r<0.99$ 标记为强共变。这些阈值是用于稳定生成 observation packet 的操作性规则，而非故障判据；参数敏感性说明可置于附录或补充材料。

\subsubsection{报告与结构化告警摘要}

每个确定性分析工具将输入范围、参数和结构化结果写入 observation packet 与记事板。compact alert summary 由其中稳定的观测字段直接生成，仅保留后续证据检索所需要的内容，例如异常事件的通道与计数、数据完整性提示、共变或滞后关系、极值位置以及选定快照线索。测量侧 agent 在不增加新的测量事实的前提下读取这些已保存结果，并将其按时间与通道上下文组织为 condition report，供工程师查看数据健康、事件、共变、滞后、极值和快照线索。

这种“完整报告—紧凑摘要”的区分避免了两种相反的问题：只保留自然语言叙述会使检索输入松散且难以复现；只保留单个统计量又会丢失工程师复核所需的时间与通道上下文。摘要因此是由脚本和记事板直接生成的受控输出；完整报告虽经语言模型组织，也只能呈现已保存的观察，而不是自由改写的诊断结论。

\subsection{证据侧：从结构化摘要到候选解释}

技术证据资源由核仪表与控制相关资料构成，并为每份文档保留来源 metadata，以支持段落定位和后续引用。当前资料库包含 64 份文档，其中 13 份构成以 IAEA 在线监测和仪表通道评估资料为主的权威核心，其余为 NRC、EPRI/DOE、OECD/NEA 及实验室或专题报告等同领域补充资料。在线监测与仪表通道评估的工程背景由 IAEA 指南和相关综述提供 \citep{iaea2008olm1,iaea2008olm2,hashemian2011online}。

核心资料与补充资料的区别不是“相关”与“不相关”的二元划分，而是来源策略的区别。补充资料扩大技术覆盖范围；核心资料则对应部署时预先指定的权威来源。因而，证据侧既需要检索与观察相符的段落，也需要记录所得证据是否满足既定来源策略。第四章将把这一要求作为检索治理问题进行评价。

\subsubsection{由告警字段约束检索意图}

compact alert summary 是测量侧与证据侧之间唯一的受控连接。它不把突刺、零值、共变或滞后直接命名为某种故障，而是将这些现象映射为待核查的技术主题。查询的具体措辞由作者根据摘要、语种和资料术语人工设定并复核，而非从原始波形自动生成物理诊断假设。

\begin{table}[H]
\centering
\footnotesize
\begin{tabular}{p{0.29\columnwidth} p{0.63\columnwidth}}
\toprule
摘要中的观察线索 & 检索意图与优先查找的证据 \\
\midrule
突刺或孤立零值 & 信号完整性和异常测量背景；定义、核查步骤与诊断注意事项 \\
共变或滞后 & 冗余、共享上下文与时序行为；仪表组织、在线监测与时间行为说明 \\
形态或数据健康线索 & 稳定性、噪声、可用性与监测实践；数据质量和监测指南 \\
极值或通道快照 & 运行上下文与通道间比较；运行指导和推荐检查项 \\
\bottomrule
\end{tabular}
\caption{由 observation packet 中的稳定字段形成受控检索意图。}
\label{tab:retrieval-intent}
\end{table}


该映射限制的是检索范围，而不是解释结论。例如，观察到某次零值事件时，系统可查找有关测量完整性或通道检查的资料段落；它不能据此自动声称出现了某一种仪表失效。

\subsubsection{资料处理、段落检索与候选解释}

资料库中的 PDF 在必要时先进行页面裁切和文本清理，再转换为可检索文本、分块并建立向量索引；每个 chunk 保留文档身份、来源类别、页码或段落定位等 metadata。检索时，系统首先寻找与已观察现象或相应技术机制相关的段落；随后可补充与冗余、通道组织、在线监测实践或建议检查有关的资料。这样，检索对象始终是可定位的技术段落，而不是一个脱离来源的语言模型回答。

候选解释以 evidence-linked memo 的形式组织：每一项表述均应关联到具体的观测字段和资料段落，并以引文标识其支持来源。资料段落支持的是某种解释的适用背景、检查方向或工程含义；它不自动把该解释升级为已确认的根因。由此，检索配置、查询语言、编码器、分块方式和来源策略影响的是证据如何可见、如何被引用，而不是哪一个配置能普遍替代工程判断。

\subsection{接口运行约束与责任边界}

两侧任务的连接还需要运行约束来保证可追溯性。原始路径或分析状态若可随意改变，就无法确认某个摘要来自哪条记录；未预先定义的处理动作若能被自由执行，就无法复现观测；长篇报告若直接作为查询，又容易将偶然细节混入证据检索。因此，接口将这些风险转化为明确的运行规则。

\begin{table}[!t]
\centering
\footnotesize
\begin{tabular}{p{0.29\columnwidth} p{0.63\columnwidth}}
\toprule
风险 & 运行约束与保留的关联产物 \\
\midrule
输入或状态漂移 & 仅初始化步骤接收原始路径，之后均解析同一 active record；保留 active-record 标识与 probe 日志 \\
未预先定义或格式错误的动作 & 只接受预先定义的分析工具，不符合要求的动作不执行；保留接受或拒绝记录 \\
长叙述直接驱动检索 & condition report 留给人工复核，检索仅读取 compact alert summary；保留报告—摘要链接 \\
缺乏证据的文字解释 & memo 中每项候选表述关联到检索段落和引文；形成 evidence-linked memo \\
\bottomrule
\end{tabular}
\caption{measurement-to-evidence interface 的运行约束与可追溯产物。}
\label{tab:interface-constraints}
\end{table}


测量侧 agent 只调用预先定义的分析工具，并基于记事板中已保存的观测组织 condition report；数值计算由确定性脚本完成，agent 不能任意更改分析状态或添加新的测量事实。证据侧 agent 以 compact alert summary 为受控输入，组织检索意图、资料段落与候选表述；资料来源与段落定位由检索层保留。工程师保留对运行上下文、物理机制、证据充分性和最终工程决定的判断权。接口的作用不是自动替代工程判断，而是将从原始记录到候选解释之间的测量依据与文献依据保持为可回溯、可核查的链条。
